\documentclass[12pt,a4paper]{article}

\usepackage[T1]{fontenc}
\usepackage[utf8]{inputenc}
\usepackage[margin=1in]{geometry}
\usepackage{setspace}
\usepackage{booktabs}
\usepackage{amsmath}
\usepackage{siunitx}
\usepackage{hyperref}
\usepackage{enumitem}

\setstretch{1.15}
\sisetup{group-separator={,},group-minimum-digits=4}

\title{Performance and Stress Testing Evaluation Report}
\author{Hamza Béji}
\date{May 16, 2026}

\begin{document}

\maketitle
\tableofcontents
\newpage

\section{Introduction and Objectives}

As part of the validation phase of our full-stack web application built on the MERN stack
(MongoDB, Express.js, React, Node.js), comprehensive performance and stress testing was conducted.
The main objective was to evaluate the stability, scalability, and response latency of the
Express.js backend and MongoDB layer under high concurrent traffic.

Testing was performed using Grafana k6 (v2.0.0-rc1), enabling observation of server behavior,
resource utilization, and request queue dynamics under increasing load.

\section{Testing Methodology and Environment}

Performance tests were executed locally using simulated concurrent workloads.

The primary evaluation metrics are:

\begin{itemize}[leftmargin=1.5em]
    \item \textbf{Throughput} $(T)$: number of successful HTTP requests processed per second
    \begin{equation}
        T = \frac{N_{\text{requests}}}{\Delta t}
    \end{equation}
    where $\Delta t$ is the total test duration (\SI{55}{s}).

    \item \textbf{Error Rate} $(E)$: percentage of failed HTTP responses
    \begin{equation}
        E = \frac{N_{\text{failed}}}{N_{\text{total}}} \times 100\%
    \end{equation}

    \item \textbf{Percentiles} $p(x)$: response-time threshold below which $x\%$ of observations fall.
    Specifically, $p(90)$ and $p(95)$ were monitored to assess tail latency.
\end{itemize}

The load profile followed three stages:

\begin{itemize}[leftmargin=1.5em]
    \item \textbf{Stage 1 (Ramp-up):} 0 $\rightarrow$ 100 Virtual Users (VUs) in \SI{15}{s}
    \item \textbf{Stage 2 (Stress):} 100 constant VUs for \SI{30}{s}
    \item \textbf{Stage 3 (Ramp-down):} 100 $\rightarrow$ 0 VUs in \SI{10}{s}
\end{itemize}

\section{Test Scenarios and Analytical Results}

\subsection{Scenario A: Read-Intensive Database Load (GET Request)}

This scenario evaluated system performance during standard data retrieval operations.

\begin{itemize}[leftmargin=1.5em]
    \item \textbf{Target endpoint:} \texttt{GET /api/doctor}
    \item \textbf{Generated volume:} 8,550 checks across 4,275 HTTP requests
\end{itemize}

\begin{table}[h!]
\centering
\caption{GET Request Performance Metrics}
\begin{tabular}{p{0.42\linewidth} p{0.20\linewidth} p{0.30\linewidth}}
\toprule
\textbf{Metric} & \textbf{Value} & \textbf{Status / Interpretation} \\
\midrule
Checks succeeded ratio & 100.00\% (8550/8550) & Passed \\
HTTP request failure rate $(E)$ & 0.00\% & Completely stable \\
Throughput $(T)$ & 76.84 req/s & Efficient resource handling \\
Average latency $(\mu)$ & 5.11 ms & Very fast response \\
Median latency & 4.76 ms & Consistent response time \\
Maximum latency & 40.66 ms & Within safe limits \\
95th percentile $p(95)$ & 8.77 ms & Excellent tail latency \\
\bottomrule
\end{tabular}
\end{table}

\textbf{Analysis.}
The read-intensive workload performed exceptionally well. An average latency of
$\mu = \SI{5.11}{ms}$ under a peak of 100 VUs indicates efficient routing, indexing, and likely
effective caching behavior. No observable degradation occurred during the stress window.

\subsection{Scenario B: Compute-Intensive Authentication Load (POST Request)}

This scenario simulated concurrent login requests requiring database access and
cryptographic password verification.

\begin{itemize}[leftmargin=1.5em]
    \item \textbf{Target endpoint:} \texttt{POST /api/user/login}
    \item \textbf{Generated volume:} 2,772 HTTP requests
\end{itemize}

\begin{table}[h!]
\centering
\caption{POST Request Performance Metrics}
\begin{tabular}{p{0.42\linewidth} p{0.20\linewidth} p{0.30\linewidth}}
\toprule
\textbf{Metric} & \textbf{Value} & \textbf{Status / Interpretation} \\
\midrule
HTTP request failure rate $(E)$ & 0.00\% (2772/2772) & Functional success \\
Throughput $(T)$ & 7.56 req/s & Significantly reduced throughput \\
Average latency $(\mu)$ & 6.77 s & Critical bottleneck \\
Maximum latency & 15.70 s & Severe peak delays \\
Time threshold check ($<400$ ms) & 5.25\% passed & 94.75\% failed \\
95th percentile $p(95)$ & 9.98 s & Severe queueing effects \\
\bottomrule
\end{tabular}
\end{table}

\textbf{Analysis and Bottleneck Identification.}
Although functional reliability remained intact ($E = 0.00\%$), a major architectural bottleneck
was identified. Average latency increased to $\mu = \SI{6.77}{s}$, with worst-case delays of \SI{15.70}{s}.

This degradation is primarily caused by CPU-bound password verification (\texttt{bcrypt.compare}).
Because Node.js relies on a single main event loop for JavaScript execution, concurrent
cryptographic operations can saturate processing capacity and increase event-loop wait time,
resulting in cascading request latency.

\section{Conclusion and Architectural Recommendations}

The system maintained full functional integrity (0\% error rate) under both stress scenarios.
However, to support realistic production-scale traffic, the following improvements are recommended:

\begin{enumerate}[leftmargin=1.5em]
    \item \textbf{Horizontal scaling with clustering.}
    Use PM2 cluster mode (or native Node.js clustering) to spawn multiple worker processes
    mapped to available CPU cores, mitigating single-process throughput limits.

    \item \textbf{Cryptographic workload offloading.}
    Offload bcrypt operations to Worker Threads or a dedicated authentication service to prevent
    main-thread saturation and reduce queueing delay.

    \item \textbf{Rate limiting and request queue control.}
    Apply throttling and backpressure at the API gateway to protect authentication endpoints
    from burst traffic and preserve global responsiveness.

    \item \textbf{Optional algorithm/cost-factor tuning.}
    Reevaluate bcrypt cost factor to balance security and latency objectives, and benchmark
    alternatives only if security policy permits.
\end{enumerate}

\end{document}
